% Generated from locale/tr/content/first-order-logic/tableaux/derivations.tex; do not hand-edit.


\iftag{FOL}
      {\olfileid[tr]{fol}{tab}{der}}
      {\olfileid[tr]{pl}{tab}{der}}

\olsection{\usetoken{P}{tableau}}

\begin{explain}
Bir varsayımın ne olduğunu söyledik ve çıkarım kurallarını verdik.
!!^{tableau}s bunlardan tümevarımsal olarak oluşturulur: her !!{tableau} ya
bir veya daha fazla varsayımdan oluşan tek bir daldır ya da !!a{tableau}'nun
bir dalına çıkarım kurallarından birinin uygulanmasıyla elde edilir.
\end{explain}

\begin{defn}[!!^{tableau}]
$\sFmla{S_1}{!A_1}$, \dots, $\sFmla{S_n}{!A_n}$ varsayımları için bir
!!^a{tableau} (her $S_i$, $\True$ ya da~$\False$ olmak üzere), aşağıdaki
koşulları sağlayan sonlu bir !!{signed formula}s ağacıdır:
\begin{enumerate}
\item Ağacın en üstteki $n$ !!{signed formula}s{}i, alt alta yazılmış
  $\sFmla{S_i}{!A_i}$ ifadeleridir.
\item Ağaçta varsayımlardan biri olmayan her !!{signed formula}, üstündeki
  dalda bulunan !!a{signed formula}'ya bir çıkarım kuralının doğru
  uygulanmasıyla elde edilir.
\end{enumerate}
!!a{tableau}'nun bir dalı hem $\sFmla{\True}{!A}$ hem
de~$\sFmla{\False}{!A}$ içeriyorsa ve ancak o zaman \emph{kapalı}, aksi
hâlde \emph{açık}tır. Her dalı kapalı olan !!^a{tableau}, (varsayımlar kümesi
için) \emph{kapalı bir !!{tableau}}'dur. !!a{tableau} kapalı değilse, yani
en az bir açık dal içeriyorsa \emph{açıktır}.
\end{defn}

\begin{ex}
  Her varsayım kümesi tek başına !!a{tableau}'dur, ancak genellikle kapalı
  değildir. (Açıktır ki ancak varsayımlar zaten
  $\sFmla{\True}{!A}$ ile~$\sFmla{\False}{!A}$ biçimindeki bir
  !!{signed formula}s çiftini içeriyorsa kapalıdır.)

  Bir !!a{tableau}'dan (açık ya da kapalı) içindeki
  !!a{signed formula}~$\sFmla{S}{!A}$'ya
  çıkarım kurallarından birini uygulayarak yeni ve daha büyük bir tableau
  elde edebiliriz. Kural, uygulandığı~$\sFmla{S}{!A}$ geçişini içeren her dalın sonuna
  bir veya daha fazla !!{signed formula}s ekler.

  Örneğin $\sFmla{\True}{!A \land \lnot !A}$ varsayımını ele alalım.
  Yalnızca bu varsayımdan oluşan (açık) !!{tableau} şöyledir:
  \begin{center}
    \begin{tableau}{}
      [\sFmla{\True}{\formula{A} \land \lnot \formula{A}}, just=\TAss]
    \end{tableau}{}
  \end{center}
  Varsayıma $\TRule{\True}{\land}$ kuralını uygulayarak bundan yeni bir
  !!{tableau} elde ederiz. Bu kural !!{tableau}'ya iki yeni satır,
  $\sFmla{\True}{!A}$ ve $\sFmla{\True}{\lnot !A}$ eklememize izin verir:
  \begin{center}
    \begin{tableau}{}
      [\sFmla{\True}{\formula{A} \land \lnot \formula{A}}, just=\TAss
        [\sFmla{\True}{\formula{A}}, just={\TRule{\True}{\land}[1]},
          [\sFmla{\True}{\lnot \formula{A}}, just={\TRule{\True}{\land}[1]}
          ]
        ]
      ]
    \end{tableau}{}
  \end{center}
  !!{tableau}s yazarken uyguladığımız kuralları sağ tarafta kaydederiz
  (örneğin $\TRule{\True}{\land} 1$, o satırdaki !!{signed formula}'nın
  $1$.~satırdaki !!{signed formula}'ya $\TRule{\True}{\land}$ kuralının
  uygulanmasıyla elde edildiği anlamına gelir). Bu yeni !!{tableau} artık
  ek !!{signed formula}s içerir; ancak yalnızca birine
  ($\sFmla{\True}{\lnot !A}$) kural uygulayabiliriz (burada
  $\TRule{\True}{\lnot}$ kuralını). Böylece şu kapalı !!{tableau} elde edilir:
  \begin{center}
    \begin{tableau}{}
      [\sFmla{\True}{\formula{A} \land \lnot \formula{A}}, just=\TAss
        [\sFmla{\True}{\formula{A}}, just={\TRule{\True}{\land}[1]}
          [\sFmla{\True}{\lnot \formula{A}}, just={\TRule{\True}{\land}[1]}
            [\sFmla{\False}{\formula{A}}, just={\TRule{\True}{\lnot}[3]}, close]
          ]
        ]
      ]
    \end{tableau}{}
  \end{center}
\end{ex}

